% Appendix: Prior Mismatch Sensitivity Analysis

\subsection{Prior Mismatch Sensitivity (Figures~\ref{fig:prior_mismatch_heatmap}--\ref{fig:prior_mismatch_distribution})}
\label{appendix:prior_mismatch}

The warmup ablation (Section~\ref{sec:warmup_ablation}) demonstrated
that well-calibrated priors substantially reduce regret and eliminate
catastrophic cold-start failures under stationary conditions.
All experiments in this paper use priors trained on a representative
training set drawn from the same distribution as the test traffic.
In practice, however, priors may be trained on stale
data, a different prompt mix, or an outdated model portfolio.
This experiment stress-tests warmup priors under progressively worse
calibration to quantify the cost of bad priors and identify when
tabula rasa is the safer choice.

\paragraph{Design.}
We construct five prior-quality levels forming a mismatch gradient:
\begin{enumerate}[nosep,leftmargin=*]
  \item \textbf{Well-calibrated:} priors from the full training set
    ($n{=}8{,}374$)---the shipped default.
  \item \textbf{Random-1680:} a random subsample of 1,680 training
    prompts, matching the GSM8K-only count while preserving the full
    distribution.  This is a \emph{sample-size control}: any gap
    between Random-1680 and Well-calibrated isolates covariance
    estimation noise from domain effects.
  \item \textbf{MMLU-only:} priors from the 1,855 MMLU prompts only.
    The prior preserves the correct model ranking (Gemini~$>$~Mistral~$\gg$~Llama)
    but with knowledge-domain-specific magnitudes.
  \item \textbf{GSM8K-only:} priors from the 1,680 GSM8K (math) prompts only.
    All models score ${\sim}0.95{+}$ on math, so the prior encodes
    near-zero arm differentiation.
  \item \textbf{Inverted:} priors from the full training set with
    Llama and Gemini rewards swapped, so the prior believes the
    cheapest model is best and vice versa.
\end{enumerate}
Each level is tested at three $n_\mathit{eff}$ values (10, 100, 1000)
plus a single Tabula Rasa control ($n_\mathit{eff}{=}1$, no priors),
yielding $5 \times 3 + 1 = 16$ conditions.
All warmup conditions use $\alpha{=}0.01$, $\gamma{=}0.997$ (disjoint
LinUCB); Tabula Rasa uses $\alpha{=}0.01$, $\gamma{=}0.995$
(each condition's hyperparameters are independently optimised via the
Experiment~05 Pareto knee-point sweep).
The warmup hyperparameters were \emph{not} re-optimised per
prior-quality level---matching the production scenario where prior
quality degrades after deployment without hyperparameter adjustment.
Note that the $\gamma$ difference ($0.997$ vs.\ $0.995$) means
warmup conditions have ${\sim}67\%$ longer effective memory than
Tabula Rasa.
This confound is addressed by the $n_\mathit{eff}$ dose--response
within this experiment: all warmup conditions share $\gamma{=}0.997$,
yet regret decreases monotonically with $n_\mathit{eff}$
(e.g., Well-calibrated: $n_\mathit{eff}{=}10 \to 100 \to 1{,}000$),
confirming that the prior content---not the forgetting factor---drives
the advantage.
Evaluation follows the same cumulative-regret protocol as the warmup
ablation: held-out test split ($n{=}1{,}824$), 20 seeds
(offset${=}9{,}000$, aligned for paired comparisons), unconstrained
(no budget pacer).
We deliberately restrict this experiment to the unconstrained regime
to isolate the effect of prior quality from the budget-pacer
interaction.
The warmup ablation (Section~\ref{sec:warmup_ablation}) already
characterises the pacer interaction for well-calibrated priors vs.\
tabula rasa; adding budget constraints here would conflate two
failure modes (prior mismatch and pacer-induced catastrophic
failures) without additional insight, since the prior-quality
threshold for harm is determined by the bandit's learning dynamics,
not the pacer.

\paragraph{Results.}
Figure~\ref{fig:prior_mismatch_heatmap} summarises total regret
across the quality--strength grid.
Pairwise comparisons use two complementary non-parametric tests.
The \textbf{exact binomial sign test} (two-sided,
$H_0$: $P(\text{condition wins}) = 0.5$) counts per-seed paired
wins and tests for a location shift.
The \textbf{Fisher exact test} on the $2{\times}2$ catastrophic-failure
table ($H_0$: equal failure rates; one-sided, condition $<$ baseline)
tests whether the condition reduces tail risk.
Because Tabula Rasa's per-seed regret is heavy-tailed and
multimodal ($\mathit{std}{=}82.2$; see below), the standard
assumptions of the Wilcoxon signed-rank test (symmetric paired
differences) and the paired $t$-test (normally distributed
differences) are both violated.
The sign test requires only that paired observations are independent
and makes no distributional assumptions; the Fisher test requires
only a fixed marginal total.
We complement both with win counts and catastrophic failure rates
throughout.

\emph{Directionally correct priors help at every strength.}
Well-calibrated priors reduce median total regret monotonically with
$n_\mathit{eff}$: $65.2$ vs.\ $73.4$ at $n_\mathit{eff}{=}10$
(14/20 wins, $p{=}0.12$),
$61.9$ vs.\ $73.4$ at $100$ (15/20 wins, $p{=}0.041$), and
$55.5$ vs.\ $73.4$ at $1{,}000$
(19/20 wins, $p < 10^{-4}$).
The $n_\mathit{eff}{=}10$ result is directionally positive on 14/20
seeds but does not reach significance given the sign test's limited
power at this sample size.
The sample-size control (Random-1680) performs comparably to
Well-calibrated at every $n_\mathit{eff}$---confirming that the
$n_\mathit{eff}$ scaling normalises prior strength and the benefit is
not an artifact of having more training prompts.

\emph{Domain-mismatched priors still beat cold start.}
MMLU-only priors reduce regret relative to Tabula Rasa at
$n_\mathit{eff} \geq 100$ (15--16/20 wins, $p \leq 0.041$),
despite being trained on single-domain data.
GSM8K-only priors---which encode near-zero arm differentiation---reach
significance at $n_\mathit{eff}{=}1{,}000$ (15/20 wins, $p{=}0.041$)
and are directionally positive at $100$ (14/20 wins, $p{=}0.12$).
At $n_\mathit{eff}{=}10$ both domain-specific priors are directionally
positive (13/20 and 12/20 wins, $p > 0.25$) but not significant,
reflecting the weak prior's limited initial influence.
Crucially, no domain-mismatched prior \emph{increases} regret at any
$n_\mathit{eff}$: partial knowledge is always at least as good as no
knowledge.

\emph{Inverted priors cause harm that scales with $n_\mathit{eff}$.}
This is the only condition where priors hurt.
At $n_\mathit{eff}{=}10$, the damage is modest (median $78.2$ vs.\
Tabula Rasa's $73.4$; 11/20 wins, $p{=}0.82$);
at $n_\mathit{eff}{=}100$, median regret rises to $89.9$
(7/20 wins, $p{=}0.26$);
and at $n_\mathit{eff}{=}1{,}000$, to $106.9$
(6/20 wins, $p{=}0.12$).
The sign test does not reach significance at any $n_\mathit{eff}$
because Tabula Rasa's own catastrophic seeds (${\sim}25$--$30\%$
of runs; see below) inflate the baseline's regret sufficiently to
mask the inverted condition's harm in a paired comparison.
The median comparison, however, is unambiguous: inverted
$n_\mathit{eff}{=}1{,}000$ increases median regret by $46\%$
relative to Tabula Rasa ($106.9$ vs.\ $73.4$).
The mechanism is clear: the inverted prior confidently routes traffic
to the worst model, and high $n_\mathit{eff}$ forces the router to
accumulate many online observations before overriding the prior.

\paragraph{Tabula Rasa instability.}
The mean $\pm$ SE summaries in
Figure~\ref{fig:prior_mismatch_heatmap} mask a heavy-tailed,
multimodal distribution of per-seed outcomes.
Figure~\ref{fig:prior_mismatch_distribution} makes this explicit:
Tabula Rasa's per-seed total regret spans from ${\sim}57$ to
${\sim}283$ ($\mathit{std}{=}82.2$), with 25--30\% of seeds
incurring $2{\times}$+ the median regret.
By contrast, every warmup condition---including domain-mismatched
and even inverted priors---exhibits a tight, unimodal distribution
($\mathit{std} \leq 3$ for all non-inverted conditions;
$\mathit{std} \leq 5$ even for Inverted at $n_\mathit{eff}{=}1{,}000$).

The warmup ablation (Section~\ref{sec:warmup_ablation}) reveals the
mechanism: under budget constraints, the bimodality becomes stark.
Tabula Rasa seeds split into ``good'' runs (where early exploration
happens to probe cheap arms first, yielding regret comparable to
warmup) and ``catastrophic'' runs (where early UCB-driven exploration
probes the expensive arm, the budget pacer penalises aggressively,
and the router gets stuck exploiting a suboptimal cheap arm for the
remainder of the horizon).
Under budget constraints, 40\% of Tabula Rasa seeds
catastrophically fail at moderate and loose budgets
($p_\text{Fisher}{=}0.002$), and clear bimodality appears
even at tight budgets (9/20 seeds at regret~${\sim}300$ vs.\
11/20 at~${\sim}165$).
Warmup priors eliminate this failure mode entirely: zero catastrophic
seeds across all budget levels, with
$\mathit{std} \in [4.6,\, 11.0]$.

This reframes the warmup advantage.  The sign test and the Fisher
exact test answer complementary questions: the sign test detects
location shifts (warmup consistently wins under tight budget,
$p < 10^{-3}$) while the Fisher test detects differential tail risk
(warmup eliminates catastrophic failures under moderate and loose
budgets, $p < 0.005$, even when the sign test is not significant).
The mean regret reduction reported above is a \emph{conservative}
summary that averages over a coin-flip baseline; the more
operationally relevant finding is that warmup priors
\emph{eliminate catastrophic cold-start failures}---converting a
stochastic, seed-dependent outcome into a reliable one.

\paragraph{Practical implications.}
The results establish that the prior-quality threshold for harm is
high: priors must be \emph{actively wrong} (systematically inverting
model rankings) before they degrade performance.  Domain mismatch,
reduced sample size, and even near-zero arm differentiation
(GSM8K-only) do not cross this threshold.
Beyond mean regret reduction, warmup priors deliver a qualitative
reliability guarantee: all warmup conditions produce tight, unimodal
per-seed distributions, whereas Tabula Rasa's heavy tail makes
deployment outcomes seed-dependent.
For practitioners:
\begin{itemize}[nosep,leftmargin=*]
  \item When the prior is trained on representative data from the
    deployment domain, use $n_\mathit{eff}{=}1{,}000$---this
    captures the full 24\% median regret reduction ($55.5$ vs.\
    $73.4$) and eliminates catastrophic cold-start variance.
  \item When the prior comes from a related but different domain
    (e.g., a general-purpose prior applied to a specialised
    deployment), $n_\mathit{eff}{=}100$ is a safe default that
    still reduces median regret by 16\% ($61.9$ vs.\ $73.4$)
    with negligible variance penalty.
  \item When prior quality is entirely unknown, set
    $n_\mathit{eff}{=}10$.  This limits worst-case inverted-prior
    damage to $6\%$ median increase (non-significant, $p{=}0.82$)
    while still providing a directional nudge from any
    non-adversarial prior and substantially reducing outcome
    variance.
  \item Tabula rasa ($n_\mathit{eff}{=}1$) is the appropriate
    choice only when the prior is known or suspected to be
    systematically inverted---a scenario that implies a fundamental
    error in the offline reward pipeline---and even then, operators
    should expect high seed-to-seed variance.
\end{itemize}
